\subsection{All-weight low cohomology of the first quadratic repair}

\paragraph{A-S09-20 (Proved; the exact kernel and quotient of a form-four two-cell attachment).}
Let $X$ be a nonnegative strict PD cochain $S$-algebra
with torsion-free terms, coefficient $H^0(X)=S$, and
$H^1(X)=H^2(X)=H^3(X)=0$.
Let $V$ be a free abelian parameter module with
positive weights.  For linear maps
$z:V\to Z^4(X)$ and $b:V\to X^3$ of those weights,
suppose $db=2z$.  Form the actual free PD extension
\[
 Y=X\otimes\Gamma(vV)\otimes\Lambda(uV),
 \qquad |v|=2,\quad |u|=3,
 \qquad du=z,\quad dv=2u-b.                         \tag{S09.85}
\]
Let
\[
 \zeta:(S/2S)\otimes_{\mathbb F_2}(V/2V)
                         \longrightarrow H^4(X),
                 s\otimes a\longmapsto s[z(a)].      \tag{S09.86}
\]
Then, in every weight,
\[
 \begin{split}
 H^0(Y)&=S,\\
 H^1(Y)&=H^2(Y)=0,\\
 H^3(Y)&\cong\ker\zeta,\\
 H^4(Y)&\cong\operatorname{coker}\zeta.
 \end{split}                                        \tag{S09.87}
\]
These are assertions about the entire algebra,
including products with all old coefficient and
form classes; they are not restricted to the first
weight or to an isolated arity.

\begin{proof}
Choose a basis of $V$ to write $u_i,v_i,b_i,z_i$.
Every expression below uses finitely many indices.
The free PD basis separates all indicated monomials
over the old algebra $X$.  Degrees zero and one
are unchanged.  A two-form is
$c_2+\sum_i a_i^0v_i$.  In its differential the
coefficient of $u_i$ is $2a_i^0$, so torsion-freeness
forces every $a_i^0=0$.  Closed two-forms are thus
the old closed two-forms, proving the first three
claims in (S09.87).

A three-form has the unique shape
\[
                c_3+\sum_i a_i^1v_i+\sum_i a_i^0u_i.
 \tag{S09.88}
\]
For a closed form, its $v_i$ coefficient gives
$da_i^1=0$.  Since $H^1(X)=0$, write $a_i^1=dh_i^0$.
Subtracting $d\sum_i h_i^0v_i$ removes all $v_i$
terms.  In the resulting representative the
$u_i$ coefficients are closed zero-forms, hence
coefficients $s_i\in S$.  Its remaining closedness
equation is
\[
                      dc_3+\sum_i s_i z_i=0.         \tag{S09.89}
\]
Changing a chosen $h_i^0$ by a coefficient changes
$s_i$ by a multiple of two, precisely the change
coming from $d(s_i'v_i)=2s_i'u_i-s_i'b_i$.
Equation (S09.89) therefore defines a map
$H^3(Y)\to\ker\zeta$.
It is onto, because a coefficient relation among
the $[z_i]$ has an old three-form primitive as
in (S09.89).  It is one-to-one: when every $s_i$
is even, subtract the corresponding boundaries
$d(s_i'v_i)$; the remainder is a closed old
three-form and is a boundary because $H^3(X)=0$.
This proves the degree-three formula.

A four-form has the unique shape
\[
 c_4+\sum_i a_i^2v_i+\sum_i a_i^1u_i
       +\sum_{i<j}a_{ij}^0v_iv_j
       +\sum_i a_{ii}^0\gamma_2(v_i).               \tag{S09.90}
\]
In the differential, the coefficients of the
distinct monomials $u_iv_j$ in its highest new-generator
part are $2a_{ij}^0$; for $u_iv_i$ they are
$2a_{ii}^0$.  No earlier term in (S09.90)
contributes to these monomials.  Torsion-freeness
therefore makes every quadratic coefficient zero.
The $v_i$ coefficients of the remaining differential
give $da_i^2=0$.  Write $a_i^2=dh_i^1$ using
$H^2(X)=0$, and subtract $d\sum_i h_i^1v_i$.
All $v_i$ terms disappear.  The remaining
$u_i$ coefficients are closed one-forms; write
them as $dh_i^0$ and subtract $d\sum_i h_i^0u_i$.
Every closed four-form is consequently cohomologous
to an old closed four-form.  This proves surjectivity
of $H^4(X)\to H^4(Y)$.

For its kernel, suppose an old four-form is the
differential of (S09.88).  The same removal of
$a_i^1v_i$ leaves coefficients $s_i\in S$ on
$u_i$, since the differential has no new monomials.
It is therefore $dc_3+\sum_i s_i z_i$.
Conversely every $z_i=du_i$ becomes a boundary.
The kernel is exactly the image of (S09.86),
which proves the final formula.
\end{proof}

The calculation is natural when the free parameter
module is the quotient of a relative simplicial
cell.  In a simplex, choose a free complement to
the old parameter module to obtain (S09.85).
A simplicial operator changes a new generator by
its induced linear parameter map plus an old
element.  Extracting the coefficients of $u_i$
modulo two in (S09.88) commutes with this change
in cohomology, while the four-form map is induced
by the original inclusion $X\to Y$.  The map
(S09.86) is well defined on the parameter quotient
because the old boundary values of $z$ already
have specified primitives.  Thus (S09.87) gives
isomorphisms of simplicial groups in that relative
setting as well.

Here is the concrete first application to A-S09-18--19.
For each weight $s$, the free module $Q^1(P)^s$
has normalized terms zero in degrees zero and one.
Choose a basis of $N_2(Q^1(P)^s)/2$ and the associated
map
\[
 \bigoplus K(\mathbb F_2,2)
                           \longrightarrow Q^1(P)^s/2.
 \tag{S09.91}
\]
This is an actual simplicial injection: under
normalization it is the inclusion of the chosen
degree-two basis, its differential is zero because
the degree-one term is zero, and its other source
terms are zero.  The degreewise exact Dold--Kan
correspondence transfers this injection to every
simplicial degree.  Tensoring with $S/2$ and using
the natural map (S09.80) therefore gives an
injection into $H^4(P)$ in every weight.
For an actual relative form-four attachment
realizing these classes with $m=2$, (S09.87)
then proves $H^1=H^2=H^3=0$ in all weights and
identifies its new $H^4$ with the stated quotient.
The normalized degree-two part of that quotient
vanishes.  This verifies the low-cohomology effect
of the specified independent-class step.
It supplies neither relative target primitives
for later normalized presentation relations nor
the final admissible completion; those maps are
not assumed in the calculation.

The exact certificate
\texttt{code/A\_cosimplicial\_low\_repair\_certificate.py}
checks one, two, and three copies of this attachment
over $X=\Lambda(b_3)\otimes\Gamma(z_4)$, $db=2z$.
It records every differential in degrees two through
five, verifies their square-zero compositions,
and gives integral column operations with
$H^3=(\mathbb Z/2)^{r-1}$ and $H^4=0$ for $r$
copies, as predicted by the class-map kernel.
The recorded runtime was $0.003155237005557865$ seconds;
the output is
\texttt{logs/A\_cosimplicial\_low\_repair\_certificate.json}.
